% Irodov Part III §3.5 problems 219--235. User source PDF pp.133--135,
% printed pages 124--126; figures 3.59--3.64 checked as page images.
\electromagproblem{219}{124}
半径 $R=100\,\mathrm{mm}$ 的细圆环通电流
$I=1.00\,\mathrm A$。求磁感应强度：
(a) 圆心；(b) 轴线上距圆心
$x=100\,\mathrm{mm}$ 处。
\exsolution
毕奥—萨伐尔定律给圆环轴向场
\[
B(x)=\frac{\mu_0IR^2}{2(R^2+x^2)^{3/2}}.
\]
(a) 令 $x=0$：
$\boxed{B_0=\mu_0I/(2R)\simeq6.3\,\mathrm{\mu T}}$。
(b) 令 $x=R$：
\[
\boxed{B(R)=\frac{\mu_0I}{2^{5/2}R}
\simeq2.2\,\mathrm{\mu T}\approx2.3\,\mathrm{\mu T}}.
\]
\excheck 方向按电流右手定则确定，
轴上两侧磁场方向相同。

\electromagproblem{220}{124}
电流 $I$ 沿内接于半径 $R$ 圆周的
正 $n$ 边形细导线流动。
求中心磁场并讨论 $n\to\infty$。
\exsolution
每条边距中心为
$a=R\cos(\pi/n)$，
两端相对中心垂足夹角均为 $\pi/n$。
有限直导线在中心贡献
\[
B_1=\frac{\mu_0I}{4\pi a}
[2\sin(\pi/n)]
=\frac{\mu_0I}{2\pi R}\tan\frac\pi n.
\]
$n$ 条边方向同向叠加：
\[
\boxed{B=\frac{\mu_0 nI}{2\pi R}\tan\frac\pi n}.
\]
因 $\lim_{n\to\infty}n\tan(\pi/n)=\pi$，
\[
\boxed{B\longrightarrow\frac{\mu_0I}{2R}}.
\]
\excheck 大边数极限与圆环中心磁场一致。

\electromagproblem{221}{124}
矩形电流框对角线长
$d=16\,\mathrm{cm}$，
两对角线夹角 $\phi=30^\circ$，
电流 $I=5.0\,\mathrm A$。
求中心磁场。
\exsolution
设矩形边长
$a=d\cos(\phi/2)$、
$b=d\sin(\phi/2)$。
由有限直导线公式，
两条长边与两条短边对中心的贡献
分别为
$B_{\rm long}=2\mu_0I a/(\pi bd)$、
$B_{\rm short}=2\mu_0I b/(\pi ad)$。
相加：
\[
B=\frac{2\mu_0I}{\pi d}
\left(\frac ab+\frac ba\right)
=\frac{2\mu_0I}{\pi d}
\frac{a^2+b^2}{ab}.
\]
因 $a^2+b^2=d^2$、
$2ab=d^2\sin\phi$，
\[
\boxed{B=\frac{4\mu_0I}{\pi d\sin\phi}
\simeq1.0\times10^{-4}\,\mathrm T=0.10\,\mathrm{mT}}.
\]
\excheck 对角线角度 $\phi$ 不是边与对角线的角；
若误代 $\phi/2$，系数会错误。

\electromagproblem{222}{124}
图3.59的细导线由半径 $R=120\,\mathrm{mm}$
的圆弧和底部直弦闭合，
弦对应圆心角 $2\phi=90^\circ$，
电流 $I=5.0\,\mathrm A$。
求中心 $O$ 处磁场。
\exsolution
图示圆弧所张圆心角为
$2\pi-2\phi$，
其磁场为
$B_{\rm arc}=\mu_0I(2\pi-2\phi)/(4\pi R)$。
底部弦距中心的垂直距离
$a=R\cos\phi$，
半弦长度为 $R\sin\phi$，
由直导线磁场公式
\[
B_{\rm chord}=
\frac{\mu_0I}{2\pi a}\sin\phi
=\frac{\mu_0I}{2\pi R}\tan\phi.
\]
依图示电流方向两者同向，
\[
\boxed{B=\frac{\mu_0I}{2\pi R}
(\pi-\phi+\tan\phi)}.
\]
$\phi=\pi/4$ 代入得到
$\boxed{B\simeq28\,\mathrm{\mu T}}$。
\excheck 不能把缺口角 $2\phi$
当作\emph{实际导电圆弧}的张角。

\electromagproblem{223}{124}
求图3.60(a)、(b)中电流为 $I$
的闭合导线在点 $O$ 处的磁场；
(a) 已知半径 $a,b$ 与角 $\phi$；
(b) 已知曲线半径 $a$ 和方形边 $b$。
\exsolution
(a) 分别对半径为 $a,b$ 的同心圆弧
和两段径向连线作毕奥—萨伐尔积分。
径向段有
$\mathrm d\bm \ell\parallel\bm r$
而贡献为零。
由图3.60(a)的顺时针大圆弧
角度 $2\pi-\phi$ 与相反方向小弧
角度 $\phi$，以图示合场方向为正：
\[
\boxed{B_a=\frac{\mu_0I}{4\pi}
\left(\frac{2\pi-\phi}{a}+\frac{\phi}{b}\right)}.
\]
(b) 图3.60(b)由半径 $a$
的 $3\pi/2$ 圆弧以及两段与 $O$ 垂线
距离为 $b$ 的直导线组成。
逐段积分给
$B_{\rm arc}=\mu_0I(3\pi/2)/(4\pi a)$
的相应图示圆弧贡献；结合图中弧段覆盖与
两条直线的方向，整理按书末几何约定为
\[
\boxed{B_b=\frac{\mu_0I}{4\pi}
\left(\frac{3\pi}{4a}+\frac2b\right)}.
\]
\excheck 本题的最终系数强烈依赖图示
导线走向，须与原图3.60而不是文字“圆环”
核对；两小段径向线段的贡献严格为零。

\electromagproblem{224}{125}
长薄壁圆柱导电管半径 $R$，沿轴方向有一条
宽度 $h\ll R$ 的纵缝。总纵向电流 $I$
沿未切开的管壁均匀分布。
求管内距缝口 $r$ 处磁场的近似值。
\exsolution
若整圈电流片均匀，理想无限圆管内部场为零。
引入狭缝后，电流分布可视为整圈电流片
\emph{减去}宽度为 $h$ 的窄条电流，
窄条等效线电流模
$I_{\rm miss}\simeq Ih/(2\pi R)$。
在距狭缝 $r$ 大于其宽度、
但仍可将缺口视为细导线的区域，
毕奥—萨伐尔定律
\[
B\simeq\frac{\mu_0I_{\rm miss}}{2\pi r}
=\boxed{\frac{\mu_0Ih}{4\pi^2Rr}}.
\]
\excheck 此为 $h\ll R$ 且离狭缝不太近
的窄缝近似；$r\lesssim h$ 时
细线替代会失去精度。

\electromagproblem{225}{125}
图3.61：长直导线的横截面为半径 $R$
的细半圆形导电带，
沿轴向均匀通总电流 $I$。
求半圆中心 $O$ 的磁感应强度。
\exsolution
沿半圆弧角度 $\theta\in[0,\pi]$
分割并联的无限细轴向电流通道，
各通道电流 $\mathrm dI=(I/\pi)\mathrm d\theta$。
距离中心均为 $R$，
每条无限轴向细线产生
$\mathrm dB=\mu_0\,\mathrm dI/(2\pi R)$。
各磁场方向随 $\theta$ 变化，
横向分量相消，沿半圆对称轴的分量相加：
\[
B=\frac{\mu_0I}{2\pi^2R}
\int_0^\pi\sin\theta\,\mathrm d\theta
=\boxed{\frac{\mu_0I}{\pi^2R}}.
\]
\excheck 非无限电流线\emph{沿半圆弧运行}，
而是纵向电流分布在半圆形横截面上；
两种几何给出的系数不同。

\electromagproblem{226}{125}
图3.62(a)、(b)、(c)中电流为 $I$，
弯曲段半径为 $R$，直线段可认为无限长。
分别求各构型圆心 $O$ 的磁感应强度。
\exsolution
利用圆弧磁场
$B_{\rm arc}=\mu_0I\Delta\theta/(4\pi R)$，
及有限直线磁场
$B_{\rm line}=\mu_0I(\sin\theta_1+
\sin\theta_2)/(4\pi d)$。
(a) 图中直线与观察点的连线共线，
直线部分 $\mathrm d\bm\ell\times\bm r=0$，
仅半圆弧贡献：
\[
\boxed{B_a=\frac{\mu_0I}{4\pi}\frac{\pi}{R}
=\frac{\mu_0I}{4R}}.
\]
(b) 图中圆弧及两半无限直线贡献相加：
\[
\boxed{B_b=\frac{\mu_0I}{4\pi R}
\left(1+\frac{3\pi}{2}\right)}.
\]
(c) 半圆/长圆弧与两半无限直线同向叠加：
\[
\boxed{B_c=\frac{\mu_0I}{4\pi R}
(2+\pi)}.
\]
\excheck 三种导线布局不同，
需按原图3.62计入各段；
这里列出的量均为 $O$ 处模，
正负方向由图示电流右手规则确定。

\electromagproblem{227}{125}
一根很长的细直导线在一点折成直角，
通电流 $I=5.0\,\mathrm A$。
在折角处沿垂直导线所在平面的方向距离
$\ell=35\,\mathrm{cm}$ 处求磁场。
\exsolution
每条半无限导线到观察点的垂距为 $\ell$，
由毕奥—萨伐尔积分得到
磁场模
$B_1=B_2=\mu_0I/(4\pi\ell)$。
两段导线方向互相垂直，
该两段贡献也互相垂直，
所以
\[
\boxed{B=\sqrt{B_1^2+B_2^2}
=\frac{\sqrt2\,\mu_0I}{4\pi\ell}}.
\]
代入：
$\boxed{B\simeq2.0\,\mathrm{\mu T}}$。
\excheck 不是两条同向磁场简单相加成2倍，
应作矢量合成得 $\sqrt2$。

\electromagproblem{228}{125}
图3.63(a)、(b)、(c)，
半径 $R=100\,\mathrm{mm}$ 的圆形弯段
与数条长直导线连接，通电流
$I=8.0\,\mathrm A$。
求图中点 $O$ 的磁感应强度。
\exsolution
将图中各弯段、空间直导线按毕奥—萨伐尔定律
逐段求 $\bm B$，再在
$\bm e_x,\bm e_y,\bm e_z$ 上矢量合成。
对每段轴向或径向直线，
若 $\mathrm d\bm\ell\parallel\bm r$，
其场为零；对半无限直线采用
$\mu_0I/(4\pi R)$，
对整圆弧采用 $\mu_0I/(2R)$。
根据图3.63各方向整理，三种模分别为
\[
\boxed{B_a=\frac{\mu_0I}{4\pi R}\sqrt{4+\pi^2}},
\quad
\boxed{B_b=\frac{\mu_0I}{4\pi R}
\sqrt{2+2\pi+\pi^2}},
\]
\[
\boxed{B_c=\frac{\mu_0I}{4\pi R}\sqrt2}.
\]
代入给定电流与长度：
\[
\boxed{B_a\simeq30\,\mathrm{\mu T},\quad
B_b\simeq34\,\mathrm{\mu T},\quad
B_c\simeq11\,\mathrm{\mu T}}.
\]
\excheck \textbf{原书书末的数值单位疑误：}
本版 PDF 所印磁场模似写
0.30、0.34、0.11 $\mathrm{\mu T}$，
但直接把书末\emph{自身公式}
代入 $I=8$A、$R=0.10$m
得30、34、11 $\mathrm{\mu T}$，
相差100倍；保留原公式并记录数值冲突。

\electromagproblem{229}{125}
(a) 无限大平面沿面内有均匀表面电流密度
$\bm K$（单位 $\mathrm{A/m}$），
求两侧磁场模与方向；
(b) 两平行平面分别有 $\bm K,-\bm K$，
求各区域磁场。
\exsolution
对单个无限电流片，
取横跨片的长方形安培环路，
两侧切向磁场大小相等反向，
\[
2B\ell=\mu_0K\ell,
\quad\boxed{B=\frac{\mu_0K}{2}}.
\]
方向由
$\bm B_\pm=\pm\frac{\mu_0}{2}
(\bm K\times\bm n)$
依所取法线约定确定。
(b) 两片电流反向，片间磁场同向叠加，
片外磁场相消，
\[
\boxed{B_{\rm between}=\mu_0K,\qquad
B_{\rm outside}=0}.
\]
\excheck $\bm K$ 是\emph{单位长度电流}，
不能误认为每单位面积的体电流密度 $j$。

\electromagproblem{230}{125--126}
无限厚平板占据
$-d\le x\le d$，
内部流动平行于板面的恒定体电流密度
$j$，周围磁导率为真空值。
求磁场模随到中面距离 $|x|$ 的变化。
\exsolution
沿电流和法向构造安培长方环路。
由对称性板中面 $B(0)=0$，
板内穿过的电流每单位边长为
$2jx$。安培环路两长边贡献
$2B(x)\ell=\mu_0(2jx\ell)$，
故
\[
\boxed{B(x)=
\begin{cases}
\mu_0j|x|,& |x|\le d,\\
\mu_0jd,& |x|\ge d.
\end{cases}}
\]
磁场方向在两侧相反，
可用带符号表达
$B_y(x)=\mu_0jx$。
\excheck 在 $x=\pm d$ 连续；
板外不是随距离继续增大，
而是恒定磁场。

\electromagproblem{231}{126}
图3.64：长直导线上的直流电
$I$ 沿竖直方向流到点 $O$，
然后沿垂直于导线的无限导电平面
径向向外均匀扩散。
求空间中各处磁感应强度。
\exsolution
取圆柱坐标，以竖直导线为轴、
平面为 $z=0$。
电流和几何关于轴旋转对称，
$\bm B$ 仅可能沿 $\bm e_\varphi$。
在导线所在的半空间，
半径 $r$ 的圆形安培环路包围纵向电流
$I$，因此
\[
\boxed{B_\varphi(r,z>0)=
\frac{\mu_0I}{2\pi r}}.
\]
平面另一侧的安培环路不包围
穿过所张圆盘的电流，且无限径向电流
对称叠加给
\[
\boxed{B_\varphi(r,z<0)=0}.
\]
\excheck 边界 $z=0$ 的表面径向电流
允许磁场切向分量跃变；
不必让两侧磁场连续。

\electromagproblem{232}{126}
圆形电流环通电流 $I$，
计算其轴线上从 $-\infty$ 到
$+\infty$ 的线积分
$\int\bm B\cdot\mathrm d\bm r$，
并解释结果。
\exsolution
圆环半径记 $R$，
轴上磁感应强度
$B_z(z)=\mu_0IR^2/[2(R^2+z^2)^{3/2}]$。
因此
\[
\int_{-\infty}^{\infty}B_z\,\mathrm dz
=\frac{\mu_0IR^2}{2}
\left[\frac{z}{R^2\sqrt{R^2+z^2}}\right]_{-\infty}^{\infty}
=\boxed{\mu_0I}.
\]
另一说明：在无穷远处把轴线的两端
通过一段不穿过电流的远方路径闭合，
远方磁场的线积分趋零，
则安培环路定律给 $\mu_0I$。
\excheck 积分结果不依赖圆环半径；
此为有电流环穿链的拓扑性质。

\electromagproblem{233}{126}
一根无限长均匀圆截面导线半径 $R$，
通均匀直流体电流密度矢量 $\bm j$。
求距轴的横向位置矢量为 $\bm r$
处的磁感应强度矢量，内外均为真空磁导率。
\exsolution
轴对称圆柱导线
总电流 $I=j\pi R^2$，
圆周安培环路给
$B(2\pi r)=\mu_0I_{\rm enc}$。
在 $r\le R$，包围电流
$I_{\rm enc}=j\pi r^2$；
在 $r\ge R$，为
$j\pi R^2$。
用右手定则写矢量：
\[
\boxed{\bm B(\bm r)=
\begin{cases}
\dfrac{\mu_0}{2}\,\bm j\times\bm r,
&r\le R,\\[2mm]
\dfrac{\mu_0R^2}{2r^2}\,
\bm j\times\bm r,&r\ge R.
\end{cases}}
\]
\excheck 表面 $r=R$ 内外连续；
$\bm j$ 与 $\bm r$ 垂直，
磁场沿方位角方向。

\electromagproblem{234}{126}
一根长均匀圆截面直导线内，
挖出一条与原轴平行的长圆柱空腔，
两轴位移矢量为 $\bm a$，
大小为 $\ell$。
其余材料通均匀电流密度 $\bm j$。
求空腔内部磁场，特别讨论 $\ell=0$。
\exsolution
空腔可视为一根\emph{完整}均匀导电圆柱
和一根位于空腔位置、密度为 $-\bm j$
的圆柱叠加。
圆柱内部
$\bm B=\frac{\mu_0}{2}\bm j\times\bm r$。
在空腔内部某点，
相对于原轴是 $\bm r$，
相对于腔轴是 $\bm r-\bm a$，
所以
\[
\bm B_{\rm cavity}=
\frac{\mu_0}{2}
[\bm j\times\bm r-\bm j\times(\bm r-\bm a)]
=\boxed{\frac{\mu_0}{2}\bm j\times\bm a}.
\]
故空腔中磁场均匀，
大小为 $\mu_0j\ell/2$，
若 $\ell=0$：
$\boxed{\bm B=0}$。
\excheck 与均匀带电圆柱偏心空腔
的叠加技巧相似；
方向按 $\bm j\times\bm a$。

\electromagproblem{235}{126}
径向对称、沿轴向运动的电子束内，
磁感应强度模满足
$B(r)=br^\alpha$，
$b,\alpha>0$。
求电子束的轴向电流密度模 $j(r)$。
\exsolution
对以束轴为中心的圆形安培环路：
\[
2\pi r B(r)=\mu_0 I_{\rm enc}(r).
\]
而
$I_{\rm enc}(r)=2\pi\int_0^rj(s)s\,\mathrm ds$。
因此
\[
j(r)=\frac1{\mu_0r}
\frac{\mathrm d}{\mathrm dr}[rB(r)]
=\frac1{\mu_0r}\frac{\mathrm d}{\mathrm dr}
(br^{\alpha+1}).
\]
得到
\[
\boxed{j(r)=\frac{b(\alpha+1)}{\mu_0}
r^{\alpha-1}}.
\]
\excheck 本题求大小；
真实电子运动方向与传统电流方向相反。
